% 第 11 章。数字：checks/check_ch11.py；引文：data/snapshots/medieval_europe_facts.csv（逐字核对见
% data/fetch/verify_web_quotes.py）；工资、人口与地租：clark2015_england_decadal、boe_millennium_pop_england；
% 里弗尔的价值：wailly1857_livre_tournois；铸币税：rvw1996_mint_revenue；图：figsrc/fig_ch11.py。
\chapter{中世纪欧洲：庄园、黑死病与货币贬值}\label{ch:medieval}

\begin{chapterguide}
\guidequestion{一场夺走大量人口的瘟疫，为什么让幸存者的实际工资上升了？君主为什么一再让铸币贬值？}
\begin{guidegoals}
\item 用人口—土地—工资的马尔萨斯关系分析黑死病（1348—1350）前后的实际工资与地租；
\item 解释 1351 年《劳工法令》为什么压不住工资；
\item 分析铸币贬值作为财政手段的收益与代价；
\item 概括奥雷姆《论货币》对君主操纵货币的批评。
\end{guidegoals}
\guideprereq{\chref{ch:greekroman}\enspace\nameref{ch:greekroman}；\chref{ch:prices}\enspace\nameref{ch:prices}}
\end{chapterguide}

设想一个村庄一夜之间少了将近一半的人。田地、房屋、犁和牲口都还在，只是种地的人少了。留下来的人会变穷还是变富？1348—1349 年的英格兰几乎就是这样一个实验。按 Broadberry 等人的估计，英格兰的人口在三年里减少了 46\%\cite{broadberry2015}；到 15 世纪初，建筑小工与农工的实际日工资已是瘟疫前的 1.8—2.3 倍，实际地租却几乎没有下降\cite{clark2010}。同一时期，英吉利海峡对岸的法国国王一次又一次改变铸币：1351—1360 年的十年里，法国记账货币里弗尔的价值变动了 71 次\cite{wailly1857}。

这一章讨论这两件事：一场改变劳动与土地相对数量的\textbf{实物冲击}，和一连串改变计价尺度的\textbf{货币冲击}。分析前者用马尔萨斯模型（\secref{sec:medieval-malthus}）：人口减少，每个劳动者分到的土地多了，工资上升、地租下降。分析后者用铸币厂的账目（\secref{sec:medieval-debasement}）：君主每改一次铸币，就能从送到铸币厂的金银里多扣下一份。两件事之间还有两条线索。英格兰国王试图用法律把工资压回瘟疫前（\secref{sec:medieval-statute}）；法国的一位神学家奥雷姆则论证，君主改变货币是一种比高利贷还坏的征税（\secref{sec:medieval-oresme}）。最后，\secref{sec:medieval-test} 用数据检验模型：它在方向上对了，在幅度上差得很远。

\section{当时的世界}\label{sec:medieval-world}

1300 年前后的西欧以农业为主，大多数人以种地为生。按 Broadberry 等人的估计，1300 年英格兰约有 473 万人\cite{broadberry2015}。农村的土地多按\textbf{庄园}（manor）组织：领主保留一部分土地自己经营，称为\textbf{自营地}（demesne）；其余的土地由佃户耕种，佃户向领主交地租、服劳役。英格兰农民的标准份地叫\textbf{维尔盖特}（virgate）。

\subsubsection{一个庄园的账}
这一章的史料有一部分来自庄园与王室的档案。Bland、Brown 与 Tawney 在 1914 年把其中许多份译成英文出版\cite{bbt1914}，本章引用的庄园清册、瘟疫前后的账目和劳工法的执行记录都出自这部文献集。最早的一份是 1306—1307 年埃塞克斯郡哈弗林（Havering）王室庄园的清册（extent）。陪审员宣誓报告，国王在那里有自营耕地 223½ 英亩，“每英亩一年值 6 便士”，合计 111 先令 9 便士。庄园里共有 40 个维尔盖特，每年交定额地租（rent of assize）47 镑 19 先令 5½ 便士；此外，这 40 个维尔盖特的劳役“每年值 14 镑”\cite{bbt1914}。

清册把劳役逐项折成了钱。每个维尔盖特冬季要犁 4 英亩地，“每英亩的犁地值 4 便士”；耙地、脱粒、收割、锄草、运粮、运干草、施肥也一项项标了价。各项相加为 6 先令 2 便士，再加“法定增额”（lawful increment）10 便士，每个维尔盖特的劳役“一年共值 7 先令”。清册还注明，这些劳役都由佃户“自负费用”完成\cite{bbt1914}。

\begin{example}{哈弗林庄园的负担}\label{ex:medieval-manor}
用上面的数字回答：（1）40 个维尔盖特的劳役合计是不是 14 镑？（2）定额地租是劳役折价的几倍？劳役占佃户这两项负担的多少？（3）犁一英亩地的劳役，折价相当于一英亩自营地一年价值的几成？
\solution 英格兰的记账单位是 1 镑 = 20 先令 = 240 便士。（1）每个维尔盖特 6 先令 2 便士加 10 便士，即 $74+10=84$ 便士，正好 7 先令；$40\times84=3360$ 便士，正是 14 镑。（2）定额地租 47 镑 19 先令 5½ 便士合 $11\,513.5$ 便士，是劳役折价的 $11\,513.5/3360\approx3.4$ 倍；劳役占两者之和的 $3360/(3360+11\,513.5)\approx22.6\%$。（3）$4/6\approx66.7\%$。两个数字都是清册的估价，不一定是市场价格，但它们说明，1306 年的庄园连劳役也按钱来记。
\end{example}

这份清册说明了两点。第一，早在黑死病之前，佃户的负担已经主要用钱来计算和缴纳：定额地租是劳役折价的 3.4 倍。第二，劳役本身也有一个记在账上的价格。领主因此可以比较：让佃户服劳役，还是收钱再雇人。比较的结果取决于雇工的工资，而工资在半个世纪后将发生剧变。

\subsubsection{瘟疫之前：人越多，工资越低}
人口在 13 世纪持续增长。Clark 整理的英格兰数据\cite{clark2010}显示，从 1250 年代到 1300 年代，英格兰人口从 384 万增至 532 万，增加了 38\%；同期建筑小工的实际日工资下降了 24\%。人口在 1310 年代达到 561 万的高峰，实际工资又低了一截；但这十年里有 1315—1317 年的大饥荒，生活费用指数比 1300 年代高 26\%，这一段的下跌不全是人口增长造成的。（Clark 的人口序列与 Broadberry 等人的不同；本章凡与 Clark 的工资对比，都用 Clark 自己的人口序列，见本章末的“深入一步”。）人越多，工资越低，这正是下一节的模型要解释的关系。

\subsubsection{记账货币与本章的数据}
这一章说到“钱”，常常指\textbf{记账货币}。英格兰按镑、先令、便士记账，法国按里弗尔（livre）、苏（sou）、德尼耶（denier）记账，进位都是 1 : 20 : 240。契约、地租和工资都写成多少镑或多少里弗尔。在法国，每种金银铸币合多少记账单位，由国王的法令规定；国王既可以改变这个作价，也可以改变铸币的重量与成色。一里弗尔究竟含多少金银，因此是国王决定的，这是\secref{sec:medieval-debasement}的主题。

除了上面的文献集，本章还用到五类材料：1351 年《劳工法令》的原文（\emph{Statutes of the Realm} 第 1 卷，1810 年刊本）\cite{sotr1810}；Clark 整理的 1209 年以来英格兰的工资、生活费用、人口与地租序列\cite{clark2010}；Broadberry 等人重建的英格兰人口\cite{broadberry2015}；法国学者德瓦伊（Natalis de Wailly，1805—1886）1857 年对里弗尔价值的逐次计算\cite{wailly1857}；Rolnick、Velde 与 Weber（1996）整理的法国与英格兰的财政收入与铸币收入\cite{rvw1996}。奥雷姆的《论货币》据沃洛夫斯基（Louis Wolowski，1810—1876）1864 年的校订本\cite{oresme1864}。

\section{马尔萨斯模型：人口、土地与工资}\label{sec:medieval-malthus}

这一节回答：人口突然减少，工资与地租会怎样变化，变化又会持续多久？我们用 Ashraf 与 Galor（2011）的马尔萨斯模型\cite{ashrafgalor2011}。这类模型得名于马尔萨斯（Thomas Robert Malthus，1766—1834）的人口论：收入提高使人口增加，人口增加又把收入压回原处。

\subsubsection{生产与要素价格}
设经济只生产一种农产品，投入劳动 $L$ 与土地 $X$，技术水平为 $A$：
\begin{keyeq}
Y=(AX)^{\alpha}L^{1-\alpha},\qquad 0<\alpha<1 .\label{eq:medieval-output}
\end{keyeq}
每个劳动者的产出是 $y=Y/L=(AX/L)^{\alpha}$：土地与技术不变时，人越多，每人的产出越少。背后是\textbf{边际报酬递减}：在固定的土地上不断加人，每多一个人带来的产出越来越少，平均产出也就随之下降。

Ashraf 与 Galor 只关心人均收入。如果土地属于领主、劳动在市场上雇用，还可以把产出分成工资与地租。在竞争的市场上，工资等于劳动的边际产出，$w=\partial Y/\partial L=(1-\alpha)Y/L$；每英亩的地租等于土地的边际产出，$s=\partial Y/\partial X=\alpha Y/X$（本章的 $s$ 指地租，与\appref{app:symbols}中表示基本盈余的 $s$ 无关）。劳动得到产出的 $1-\alpha$，土地得到 $\alpha$。设技术与土地不变，劳动从 $L_0$ 降到 $L_1$，则
\begin{keyeq}
\frac{w_1}{w_0}=\left(\frac{L_1}{L_0}\right)^{-\alpha},\qquad \frac{s_1}{s_0}=\left(\frac{L_1}{L_0}\right)^{1-\alpha} .\label{eq:medieval-factor}
\end{keyeq}
劳动减少，工资上升、地租下降；$\alpha$ 越大，即土地在生产中越重要，工资涨得越多。

\begin{intuition}[人少了，每个人分到的地多了]
一个村庄有 100 英亩地。人多的时候，同一块地上挤着更多的人，地翻得更细、草除得更勤，可是多加一个人，收成只多一点。最后一个劳动者多带来的产出，决定了雇主最多愿意付多少工资。人少了一半，每个人分到的地多了一倍，最后一个劳动者多带来的产出高了，工资随之上涨。（人多时还要去种贫瘠的边角地，人少时把它们放弃，这条渠道的方向相同；本章的模型只有一种质量的土地，没有把它单列。）反过来，地多人少，领主找不到佃户，地租就要下降。黑死病没有带走一寸土地，它改变的是劳动与土地的相对稀缺程度。
\end{intuition}

\subsubsection{人口的反应}
模型的第二部分是人口。Ashraf 与 Galor 设每个成年人把收入的固定比例 $\gamma$ 用于养育子女，养活一个孩子的成本为 $\rho$，于是每人存活的子女数为 $n_t=\gamma y_t/\rho$，下一代的劳动者为 $L_{t+1}=n_tL_t$\cite{ashrafgalor2011}。收入越高，存活的孩子越多。把它与式~\eqref{eq:medieval-output} 合起来，人均收入逐代按下面的方程变化：
\begin{keyeq}
y_{t+1}=\Big(\frac{\rho}{\gamma}\Big)^{\alpha}y_t^{\,1-\alpha},\qquad \bar y=\frac{\rho}{\gamma},\qquad \frac{y_{t+1}}{\bar y}=\Big(\frac{y_t}{\bar y}\Big)^{1-\alpha} .\label{eq:medieval-dynamics}
\end{keyeq}
这里的一期是一代人。稳态的人均收入 $\bar y$ 只取决于生育的偏好与养育的成本，与技术和土地无关：技术进步或新开垦的土地只会让人口变多，不会让人均收入长期提高。这就是马尔萨斯式经济的结论；欧洲怎样走出这个陷阱，是\chref{ch:divergence}的问题。人口按同样的规律变化：$L_{t+1}=(\gamma/\rho)(AX)^{\alpha}L_t^{1-\alpha}$，Ashraf 与 Galor 把右边记作 $\varphi(L_t)$；以稳态人口 $\bar L$ 为单位，$L_{t+1}/\bar L=(L_t/\bar L)^{1-\alpha}$。

Ashraf 与 Galor 正好用黑死病作例子：像黑死病（1348—1350）那样的瘟疫，“会暂时减少人口，同时暂时提高人均收入”\cite{ashrafgalor2011}。“暂时”是关键：收入高了，存活的孩子多了，人口逐代回升，收入逐代回落。式~\eqref{eq:medieval-dynamics} 说，人均收入对稳态的偏离（取对数）每过一代只剩上一代的 $1-\alpha$ 倍。

\subsubsection{$\alpha$ 有多大}
要得到数字，需要知道土地的份额 $\alpha$。Clark 估计，1300—1349 年英格兰农业的产值中，“工资占 50\%，土地占 34\%”\cite{clark1999wp}，其余 16\% 归资本。模型里只有劳动与土地，所以要把资本处理掉。Clark 的做法是假定“资本的回报率在 1349 年之后保持不变”\cite{clark1999wp}，资本随产出同比例调整。在这个假定下，两要素模型中的土地份额是 $\alpha=0.34/(1-0.16)\approx0.405$，推导见下面的“深入一步”。

\begin{deeper}[资本可以调整时的土地份额]
设产出为 $Y=BK^aL^bX^c$，$a+b+c=1$，$a$、$b$、$c$ 分别是资本、劳动、土地的份额，$B$ 是技术。如果资本的回报率 $r$ 固定，资本会调整到它的边际产出等于 $r$：$aY/K=r$，即 $K=aY/r$。代回生产函数，
\[
Y=B\Big(\frac{aY}{r}\Big)^{a}L^{b}X^{c}\quad\Longrightarrow\quad Y=\Big[B\Big(\frac{a}{r}\Big)^{a}\Big]^{\frac{1}{1-a}}L^{\frac{b}{1-a}}X^{\frac{c}{1-a}} .
\]
因为 $\frac{b}{1-a}+\frac{c}{1-a}=1$，这是一个只有劳动与土地、规模报酬不变的生产函数，与式~\eqref{eq:medieval-output} 同形，土地的份额为 $\alpha=c/(1-a)$。用 Clark 的份额，$\alpha=0.34/0.84\approx0.405$，劳动的份额 $1-\alpha\approx0.595$，正好等于 $0.50/0.84$。如果资本的数量反而固定不变，工资对劳动减少的反应会更大：每个劳动者分到的资本也变多了。
\end{deeper}

\begin{example}{黑死病的模型预测}\label{ex:medieval-plague}
取 $\alpha=0.405$。（1）Clark 的工作论文估计，1400 年的劳动力是 1300—1349 年的 0.44\cite{clark1999wp}。按式~\eqref{eq:medieval-factor}，实际工资与地租各变化多少？（2）按 Clark 后来整理的人口序列，1400 年代的人口为 264 万，1300—1349 年平均为 501 万，比值约 0.53。预测又是多少？（3）如果此后不再有瘟疫，按式~\eqref{eq:medieval-dynamics}，（2）中人均收入的增幅一代之后、四代之后各剩多少？多少代之后，增幅（取对数）减半？
\solution （1）$0.44^{-0.405}\approx1.39$，$0.44^{0.595}\approx0.61$：工资上涨 39\%，地租下降 39\%。Clark 自己的计算是，人口减少一半以上，到 1400 年“只带来约 40\% 的实际工资上涨和 40\% 的实际地租下降”；更准确地说，“如果工资上涨 40\%，地租应当下降 39\%”\cite{clark1999wp}。（2）$0.53^{-0.405}\approx1.30$，$0.53^{0.595}\approx0.68$：工资上涨 30\%，地租下降 32\%。（3）设 $x_t=y_t/\bar y$，$x_0=1.30$，$x_{t+1}=x_t^{0.595}$：一代之后 $x_1\approx1.17$，四代之后约 1.03，只高出稳态 3\%。增幅的对数每代乘以 0.595，减半需要 $\ln0.5/\ln0.595\approx1.3$ 代。
\end{example}

图~\ref{fig:malthus_dynamics} 画出了这个过程。人口从稳态的 0.53 出发，每一代都比上一代多，但增加得越来越少；人均收入则从高出稳态 30\% 的位置逐代回落。模型给出了三个可以检验的预测：工资上涨、地租下降，两者的幅度由 $\alpha$ 决定；工资的上涨是暂时的，几代之内消失；人口会逐步回到原来的水平。

\fig{malthus_dynamics}{马尔萨斯模型中瘟疫之后的人口与人均收入，$\alpha=0.405$。(a) 看折线怎样在曲线与 45° 线之间来回：人口从稳态的 0.53 出发，每一代先竖直走到曲线（下一代的人口），再水平回到 45° 线，一步步回到稳态。(b) 看两条线怎样相向收敛：瘟疫后的第一代人均收入高出稳态 30\%，此后逐代回落，四代之后只高出约 3\%。数据：人口降幅按 Clark 的人口序列（1400 年代对 1300—1349 年平均）；份额按 Clark 对 1300—1349 年的估计；模型见 Ashraf 与 Galor（2011）。}{fig:malthus_dynamics}

模型里的瘟疫只发生一次。现实中，瘟疫在此后的一百多年里反复出现。Clark 认为，1210—1500 年英格兰工资与人口的起伏，可能完全来自人口与疾病，而与技术无关：“先是一段高生育率或低死亡率的时期，然后是一段低生育率或高死亡率的时期”\cite{clark1999wp}。在 Ashraf 与 Galor 的框架里，反复的瘟疫可以理解为养活一个能长大成人的孩子的成本 $\rho$ 上升了；由式~\eqref{eq:medieval-dynamics}，$\rho$ 上升会提高稳态收入 $\bar y=\rho/\gamma$。这不是 Ashraf 与 Galor 的结论，但与 Voigtländer 与 Voth（2013）的解释思路相近。他们认为，黑死病这样的大冲击“可以触发向一个人均收入更高的新稳态的过渡”：高工资使城市扩大，城市不卫生，战争与贸易又传播疾病，死亡率因此长期偏高，“高工资本身就这样减轻了人口压力”\cite{vv2013}。\secref{sec:medieval-test} 会用这个思路来理解人口为什么迟迟没有恢复。

\section{黑死病}\label{sec:medieval-plague}

这一节回答：瘟疫带走了多少人，当时的人又看到了什么？

英格兰的人口没有逐年的统计。Broadberry、Campbell、Klein、Overton 与 van Leeuwen（2015）重建了若干基准年的人口，基准年之间用插值\cite{broadberry2015}。按他们的序列，英格兰人口 1348 年为 481 万，1351 年为 260 万，减少了 46\%；1377 年的基准估计为 250 万，1400 年为 208 万，1450 年为 190 万，只有 1348 年的 40\%。1348 到 1351 年的下降是这个序列对 1348—1349 年瘟疫的表示，不是逐年的估计。Humphries 与 Weisdorf 说，黑死病让不列颠“几乎一半的人口”死去\cite{hw2019}。Clark 的人口序列水平不同，1340 年代为 448 万，1400 年代为 264 万，但它同样显示人口在一个世纪里没有恢复（图~\ref{fig:plague_wages}(a)）。

当时的文书记下了瘟疫的直接后果。1350 年，约克教区塞尔比（Selby）修道院的院长写信给国王，委托两名代理人替他出席议会。他给出的理由是，本院“那些较审慎、较强壮的弟兄，掌管本院事务的人，已因瘟疫而走了世人都要走的路”，而本院“在地租的减少和粮食与其他食物的匮乏中”正遭受不应有的灾难\cite{bbt1914}。

“地租的减少”可以在账目里看到。南弗里斯（South Frith）是一处有铁场的领地，收租人托马斯·贾德（Thomas Judde）每年要报告收不上来的地租（defect of rent）。1345—1346 年，这一项合计 15 先令 4 便士，其中最大的一笔是 40 英亩地，“因为没有佃户而荒着”，每年 13 先令 4 便士。1349—1350 年，同一项合计 119 先令 3¼ 便士，账上写明，其中 103 先令 11¼ 便士“是因瘟疫而新出现的短缺”\cite{bbt1914}。收不上来的地租是瘟疫前的 7.8 倍。

一封信、一份账目不能代表整个英格兰，但它们与\secref{sec:medieval-malthus}的模型说的是同一件事：人少了，土地找不到耕种的人，领主的地租收入下降。劳动这一边的变化更快。瘟疫还没有过去，国王就开始立法压工资了。

\section{劳工法令：把工资压回瘟疫前}\label{sec:medieval-statute}

这一节回答：用法律把工资定在瘟疫前的水平，为什么压不住？

\subsubsection{1349 年条例与 1351 年法令}
1349 年 6 月 18 日，爱德华三世（Edward III，1327—1377 年在位）以令状颁布《劳工条例》（Ordinance of Labourers）。序言说：“由于大部分人民，特别是工人和仆役，已在这场瘟疫中死去，有些人看到领主的需要和仆役的稀缺，除非得到过高的工资就不肯服务。”\cite{bbt1914}条例规定，王国内“六十岁以下”、身体健全、没有其他生计的男女，都必须受雇于需要他的人，只能领取“在我们在位的第二十年，或此前五六个普通年份”惯常的工资。爱德华三世在位的第二十年是 1346—1347 年。付出、许诺、索取或收受更高的工资，罚以“所付、所诺、所索或所收数额的两倍”；庄园或市镇的领主违反条例，罚三倍；出售食物的人也必须“以合理的价格出售”\cite{bbt1914}。

议会在 1351 年 2 月通过了《劳工法令》（Statute of Labourers）\cite{sotr1810}。法令的序言抱怨，仆役不顾条例，“除非得到第二十年及以前惯常所得的两倍或三倍的衣食与工资”就不肯服务。法令把条例的原则写得更具体。农业仆役要“按整年或其他通常的期限受雇，而不是按日”，工资以第二十年“或此前四年”为准；晒草、锄草的季节每天只付 1 便士，割草按英亩或按日都是 5 便士，八月第一周的收割工 2 便士、第二周 3 便士，都“不供饭食饮料”。工匠的日工资也逐项规定（第 3 章）：木匠师傅 3 便士，其他木匠 2 便士；石匠师傅 4 便士，其他石匠 3 便士，他们的帮工 1½ 便士；瓦匠 3 便士，瓦匠的小工 1½ 便士。这些同样不供饭食，是“从复活节到米迦勒节”的夏季工资，此后按专门委派的法官的裁量降低。拒绝宣誓遵守法令或违背誓言的人，可以被戴上足枷“三天或更久”，或送进监狱；被查办之后又违背誓言的，监禁 40 天，再犯监禁三个月，“每犯一次，刑罚加倍”\cite{sotr1810}。

表~\ref{tab:medieval-statute} 把这些上限与 Clark 整理的名义日工资放在一起。法令的上限大致就是瘟疫前的工资：1340 年代建筑工匠的平均日工资是 2.89 便士，比木匠师傅的上限低 4\%；建筑小工是 1.61 便士，比帮工的上限高 7\%。1350 年代，工匠的工资已经比 3 便士的上限高 35\%，小工比 1½ 便士的上限高 52\%。Clark 的序列是十年的平均日工资，不按法令的工种与季节区分，与上限只能粗略对照。

\begin{table}[htbp]\centering\small
\caption{1351 年《劳工法令》的工资上限与市场日工资（名义，便士/日）}\label{tab:medieval-statute}
\begin{tabularx}{\linewidth}{@{}>{\raggedright\arraybackslash}p{0.29\linewidth} c >{\raggedright\arraybackslash}p{0.12\linewidth} c c c@{}}
\toprule
法令中的工种 & 上限 & Clark 序列 & 1340 年代 & 1350 年代 & 1360 年代 \\\midrule
木匠师傅；其他木匠 & 3；2 & \multirow{3}{=}{建筑工匠} & \multirow{3}{*}{2.89} & \multirow{3}{*}{4.06} & \multirow{3}{*}{4.45} \\
石匠师傅；其他石匠 & 4；3 & & & & \\
瓦匠 & 3 & & & & \\
石匠的帮工、瓦匠的小工 & 1½ & 建筑小工 & 1.61 & 2.28 & 2.75 \\
晒草锄草；收割；割草 & 1；2—3；5 & 农工 & 1.49 & 2.80 & 2.88 \\
\bottomrule
\end{tabularx}
\datasource{上限：《劳工法令》第 1、3 章（\emph{Statutes of the Realm} 第 1 卷，第 311—312 页），都不供饭食，工匠为复活节至米迦勒节的夏季工资。市场工资：Clark 整理的十年平均名义日工资（\texttt{data/snapshots/clark2015\_england\_decadal.csv}），不按法令的工种与季节区分，只能与上限粗略对照。计算见 \texttt{checks/check\_ch11.py}。}
\end{table}

\subsubsection{为什么压不住}
用\secref{sec:medieval-malthus}的模型看，法令是一种\textbf{价格上限}。劳动需求曲线由式~\eqref{eq:medieval-factor} 给出：$w=w_0(L/L_0)^{-\alpha}$，它经过瘟疫前的点 $(L_0,w_0)$。法令把工资压在 $w_0$，雇主在这个工资下想雇的人数正好是瘟疫前的 $L_0$；而劳动者只剩 $L_1$。缺口 $L_0-L_1$ 是法令自己造成的。

\begin{example}{压在旧工资上的上限}\label{ex:medieval-ceiling}
法令颁布于 1349—1351 年，按 Clark 的人口序列，1350 年代的人口是 1300—1349 年平均的 0.71；1400 年代的 0.53（例~\ref{ex:medieval-plague}）要再过半个世纪才出现。设劳动从 $L_0=1$ 降到 $L_1=0.71$，$\alpha=0.405$，法令把实际工资限在瘟疫前的 $w_0=1$。（1）没有法令时，市场工资是多少？（2）在法定工资下，雇主想雇多少人？缺口有多大？（3）设想一个人手不够的雇主，他最多愿意为第 0.71 个劳动者付多少？
\solution （1）由式~\eqref{eq:medieval-factor}，$w_1=0.71^{-0.405}\approx1.15$，比法定工资高 15\%。（2）在 $w_0=1$ 时，需求曲线给出 $L=1$：雇主想雇的人数与瘟疫前一样多，而只有 0.71 个人可雇，缺口为 $1-0.71=0.29$，将近瘟疫前劳动力的三成。（3）第 0.71 个劳动者的边际产出就是市场工资 1.15。只要多付的部分不超过 15\%，雇主雇他就有利可图：违法多付，或者用法令管不到的方式（饭食、衣物、住处）补足差额。
\end{example}

图~\ref{fig:wage_ceiling}(a) 画出了这个缺口，(b) 画出了市场上的名义日工资。法令要发挥作用，就得阻止每一笔这样有利可图的交易，而交易的双方都想成交。

\fig{wage_ceiling}{法定工资上限与市场日工资。(a) 看两条竖线之间的箭头：上限压在瘟疫前的工资，雇主想雇的人数仍是瘟疫前的劳动力，而 1350 年代的劳动者只剩 0.71，缺口为 0.29；没有上限时，工资会涨到原来的 1.15 倍。(b) 看折线何时越过表示上限的水平点线：1340 年代的日工资在上限附近，1350 年代起明显高于上限。数据：(a) 按式~\eqref{eq:medieval-factor}，$\alpha=0.405$，劳动降幅按 Clark 的人口序列（1350 年代）；(b) Clark 整理的建筑工匠与小工的十年平均日工资，上限见表~\ref{tab:medieval-statute}。}{fig:wage_ceiling}

\subsubsection{执行的记录}
法令并非一纸空文。1351 年，埃塞克斯郡切姆斯福德（Chelmsford）百户区的陪审员向劳工法官报告：仆役威廉·登克（William Denk）从雇主杰弗里·史密斯那里每年拿 20 先令，还在雇主家吃饭，“而他按法令只该拿 8 先令”，是法定数额的 2.5 倍；两名裁缝等三人收取的工钱“是他们以往所收的三倍”；大巴道（Great Baddow）的理查德·史密斯“通常收取以往两倍的工钱”\cite{bbt1914}。同一年，莱尔德拉黑（Layer de la Hay）村有 18 人因“工资超额”被罚款，从 6 便士到 40 便士不等，合计 16 先令 10 便士\cite{bbt1914}。

超额工资与罚款流向哪里，透露了执行的动力。1350 年的一份令状要求，工人超额所得的工资要“立即交给他们所属各村镇的分税官与分收税官，用于缴纳上述的十分之一与十五分之一税”，也就是用来抵扣本村应交的议会税；这份令状还规定，罚金本身仍归国王\cite{bbt1914}。到了下面这份税册，抵税的已经明确是“工人的罚款”。在埃塞克斯的温斯特里（Winstree）百户区，1351—1352 年全百户区交税 19 镑 10 先令 2 便士，其中来自居民的十五分之一税只有 38 先令 11 便士，来自工人罚款的有 17 镑 11 先令 3 便士，占 90\%。文献集的编者注明，这个百户区有一半村镇的税额完全由罚款抵付，而法令在埃塞克斯“执行得比别处更严”\cite{bbt1914}。村民举报多拿工钱的工人，等于给自己减税。

即便如此，工资还是涨了上去。Clark 认为，法令“大概被普遍规避，但规避的方式可能是在工作时给工人额外的饭食”，或者在账上把工资记得比实际付的低\cite{clark1999wp}。前一种正是例~\ref{ex:medieval-ceiling} 所说的补足差额的办法；两种办法都不会出现在以现金计的工资序列里。三十多年后，理查二世（Richard II，1377—1399 年在位）1388 年的法令重申，以往关于工匠、劳工、仆役与食物商贩的法令“都要严格遵守、切实执行”，还规定仆役或劳工，“不论男女”，雇期结束时不得离开所在的百户区到别处受雇，除非持有写明离开理由与返回时间的盖印文书\cite{bbt1914}。对流动的限制并不是新东西：1349 年的条例已禁止受雇的人在约定的雇期结束前离开雇主\cite{bbt1914}；1351 年的法令又规定，冬天住在哪个村镇，夏天只要本地有活可干，就不得离开它去别处受雇\cite{sotr1810}。1388 年的法令把这些规定收紧成一种通行证制度：工资压不住，就设法把人留在原地。

\begin{compare}{法定工资上限与市场工资}{法定工资上限（1349、1351 年）}{市场工资}
\cmprow{工资}{压在 1346—1347 年的水平}{随劳动的稀缺上涨，模型预测约 30\%}
\cmprow{谁得益}{雇主，特别是需要大量雇工的领主}{幸存的劳动者}
\cmprow{劳动怎样分配}{缺口靠违法加价、饭食补贴与强制受雇分配}{价格分配：谁出的价高，劳动就流向谁}
\cmprow{执行的代价}{法官、陪审员、罚款与监禁；罚款又成了地方税的来源}{无}
\cmprow{结果}{工资仍然上涨，立法对劳动者流动的限制越收越紧}{——}
\cmpsum{上限没有改变劳动的稀缺，只改变了稀缺怎样被分配、多付的钱以什么形式出现。}
\end{compare}

价格上限在劳动以外也屡见不鲜。罗马皇帝戴克里先（Diocletian，284—305 年在位）在 301 年颁布的限价令是一个更早、规模更大的例子（\chref{ch:greekroman}），它面对的价格上涨与货币的贬值分不开。这把我们带到本章的第二个主题。

\section{铸币贬值：君主的收入}\label{sec:medieval-debasement}

这一节回答：君主改变铸币能得到什么，谁为此付出代价？

\subsubsection{铸币厂的账}
中世纪的铸币厂是一项公共服务，也是一个税卡。任何人都可以把金银送进铸币厂，换回铸好的钱币。Rolnick、Velde 与 Weber 说，送金属的人“得到的纯金属总是比送来的少。被扣下的部分，一部分用来支付铸造成本，称为\textbf{铸造费}（brassage）；其余作为利润或税收交给君主，称为\textbf{铸币税}（seigniorage）”\cite{rvw1996}。设一马克纯金属可以铸成价值 $ME$ 个记账单位的钱币（mint equivalent），铸币厂付给送金属的人 $MP$ 个记账单位（mint price）。按他们的定义，“毛铸币税就是铸币厂扣下的差额 $ME-MP$，毛铸币税率是 $1-MP/ME$”：
\begin{keyeq}
\text{毛铸币税率}=\frac{ME-MP}{ME}=1-\frac{MP}{ME} .\label{eq:medieval-seigniorage}
\end{keyeq}
毛铸币税减去铸造成本，才是君主的净收入。按 Rolnick、Velde 与 Weber 的定义，钱币所含的纯金属减少，叫\textbf{贬值}（debasement）；增加，叫\textbf{加强}（reinforcement）\cite{rvw1996}。国王也可以不动钱币，只提高它的法定作价。两种做法都使每一记账单位所含的金属变少。贬值之后，一马克金属可以铸出更多的记账单位，$ME$ 上升；只要 $MP$ 提高得少一些，铸币税率就上升。

法国与英格兰的差别很大。Rolnick、Velde 与 Weber 统计，1285—1490 年法国的银币贬值了 123 次，其中 112 次使银含量减少 5\% 以上，最大的一次减少了 50\%；英镑则“一次从未损失超过 20\%”，直到 1542—1551 年，亨利八世（Henry VIII，1509—1547 年在位）与爱德华六世（Edward VI，1547—1553 年在位）时期的十次贬值使英镑损失了 83\% 的银含量，毛铸币税率从 2\% 升到 57\%\cite{rvw1996}。1354—1490 年的法国，正常年份的毛铸币税率银币为 7.5\%、金币为 2.0\%，贬值的年份平均为 21.7\% 与 4.3\%；1401 年的铸造成本，银币约 3\%，金币约 0.5\%。贬值还会吸引大量金属：1412 年贬值之后，铸币量增加了将近 130 倍\cite{rvw1996}。

\begin{example}{一次贬值能多收多少}\label{ex:medieval-seigniorage}
用法国 1354—1490 年银币的平均数字：正常年份毛铸币税率 7.5\%，贬值年份 21.7\%，铸造成本 3\%。（1）两种年份里，送来一马克银，铸币厂付出的 $MP$ 是 $ME$ 的多少？（2）设贬值的那一年送来的银是正常年份的 2 倍（教学假设，1412 年的实际倍数远大于此）。贬值年份的毛铸币税收入（按白银计）是正常年份的几倍？扣除成本之后呢？
\solution （1）由式~\eqref{eq:medieval-seigniorage}，$MP/ME=1-7.5\%=92.5\%$ 与 $1-21.7\%=78.3\%$。（2）毛收入之比为 $2\times21.7/7.5\approx5.8$。净铸币税率分别为 $7.5\%-3\%=4.5\%$ 与 $21.7\%-3\%=18.7\%$，净收入之比为 $2\times18.7/4.5\approx8.3$。成本是固定的一块，税率提高时，净收入增加得比毛收入更快。
\end{example}

表~\ref{tab:medieval-mint} 是两位作者从法国与英格兰的财政史研究中汇总的铸币收入，他们自己说这些数据“非常零散”\cite{rvw1996}。规律仍然清楚：法国有铸币收入数据的五个贬值年份，铸币收入至少占全部收入的 50\%，1419—1420 年高达 91\%；其余年份最多 4\%。按他们的概括，正常年份的铸币税“是微不足道的收入来源，通常在 5\% 以下；但在贬值的年份，铸币税可以占收入的 50\% 以上，如 1299、1327 与 1349 年”\cite{rvw1996}。

\begin{table}[htbp]\centering\small
\caption{法国与英格兰的年度财政收入与铸币收入}\label{tab:medieval-mint}
\begin{tabularx}{\linewidth}{@{}l l Y Y r@{}}
\toprule
& 时期 & 全部收入（千） & 铸币收入（千） & 铸币占比 \\\midrule
法国（里弗尔） & 1299* & 1965.0 & 978.7 & 50\% \\
& 1322 & 477.4 & 0.5 & 0\% \\
& 1327*† & 1254.5 & 786.7 & 63\% \\
& 1329§ & 1150.0 & 41.6 & 4\% \\
& 1349*§ & 1954.4 & 1380.0 & 71\% \\
& 1388—1389 & 2500.0 & 17.9 & 1\% \\
& 1418* & 674.6 & 513.1 & 76\% \\
& 1419—1420* & 1151.9 & 1053.3 & 91\% \\\midrule
英格兰（镑） & 1463—1466* & （30） & 5.2 & 17\% \\
& 1470—1483 & 27 & 0.5 & 2\% \\
& 1505—1509 & 142 & 0.1 & 0\% \\
& 1530 & 100 & 0 & 0\% \\
& 1544—1547* & 570 & 150.0 & 26\% \\
\bottomrule
\end{tabularx}
\datasource{Rolnick、Velde 与 Weber（1996）表 2、表 3，按当时的记账货币计。* 当年或上一年（英格兰：这一时期）发生贬值；† 上半年、§ 下半年，按年率计；1544—1547 年英格兰的铸币收入是 1544—1551 年的年平均，括号里的 30 照原表。表中只列两项都有数字的年份，全部 26 行见 \texttt{data/snapshots/rvw1996\_mint\_revenue.csv}，计算见 \texttt{checks/check\_ch11.py}。}
\end{table}

\subsubsection{里弗尔的价值}
贬值对持币人和契约双方意味着什么，要看一里弗尔究竟值多少金银。德瓦伊根据历代法令给各种金银铸币规定的作价与铸币的成色重量，计算了每一次变动之后一里弗尔所含的金银，折成 1857 年的法郎（一法郎是 5 克九成成色的银）\cite{wailly1857}。这是\textbf{按官定作价}算出的价值，不是市场上的兑换价格；而且有些年份找不到法令原文，德瓦伊只能按金银价格的变动推算，表中用字母 C 标出。图~\ref{fig:livre_tournois} 画出了 1258—1360 年的结果。

\fig{livre_tournois}{1258—1360 年一图尔里弗尔合 1857 年的法郎数，按官定的铸币作价计算。看两段阴影：腓力四世在位时，里弗尔从 20 法郎左右跌到 6.15 法郎，又在 1306 年恢复到 17.63 法郎；1351—1360 年变动了 71 次，最低 3.22 法郎。1296—1305 年的 16 次变动中有 15 次是德瓦伊在缺少法令原文时的推算值。数据：德瓦伊（1857）表五，第 221—224 页；每个数值保持到下一次变动。}{fig:livre_tournois}

圣路易（Louis IX，1226—1270 年在位）时代，一里弗尔合 20.26 法郎；1278 年 12 月以后为 20.10 法郎，并一直保持到 1295 年 4 月。腓力四世（Philippe IV，1285—1314 年在位）随后开始了一连串的贬值，下面的案例记录了它的经过与代价。

\begin{case}{法国 · 1295—1306 · 腓力四世的贬值与恢复}
1295 年 4 月，一里弗尔合 16.72 法郎；1303 年 8 月只剩 6.97 法郎\cite{wailly1857}；1305 年 4 月跌到最低点 6.15 法郎，只有他即位时的 30.6\%，跌了 69\%。德瓦伊评论说，当交易得不到保障，“当债权人的贪婪或债务人的恶意可以援引法律的不公条款，任意改变契约的真实意义”时，价格不可能达到均衡\cite{wailly1857}。

1304 年前后，高级教士向国王请愿：他们希望今后“除非在绝对必要的情况下，并由国王与枢密院成员宣誓证明”，不得改变货币。国王没有接受：他既不同意免除主教的兵役，“尤其不肯作出”他们在请愿书末尾要求的承诺\cite{wailly1857}。

1306 年 10 月，国王把货币恢复到 17.63 法郎，是一年半以前的 2.87 倍。恢复同样要付代价。德瓦伊写道，恢复所需的资金，“至少有一部分是用一种不公的手段筹措的：1306 年 7 月 22 日，宣布放逐犹太人，没收他们的财产”\cite{wailly1857}。王国境内的犹太人被逐出家园，他们的财产成了恢复货币的一部分本钱。

恢复也伤害了另一批人。巴黎的小租户早已习惯租金随货币贬值而一再缩水，“直到不足原值的三分之一”，现在被要求用足值的新币付租，他们拒绝了。“一场可怕的骚乱在巴黎爆发，房屋被抢劫焚烧，国王被围困在圣殿”\cite{wailly1857}；骚乱最后被武力平息，为首的人被绞死\cite{wailly1857}。到 1313 年 9 月，里弗尔定在 18.37 法郎。德瓦伊总结说，腓力四世留下的里弗尔，“大约是他即位时价值的 10/11，但在他在位的最后十九年里经历了二十二次变动”\cite{wailly1857}。
\end{case}

\begin{example}{一份固定租金的价值}\label{ex:medieval-rent}
设巴黎一处房屋在 1295 年以前租定，年租 10 里弗尔（教学假设），契约写的是记账货币。按图~\ref{fig:livre_tournois} 的数值，这份租金在 1295 年以前、1305 年 4 月、1306 年 10 月以后各合多少法郎？谁在哪个阶段受损？
\solution 1295 年以前：$10\times20.10=201.0$ 法郎。1305 年 4 月：$10\times6.15=61.5$ 法郎，是原来的 30.6\%，房东损失了近七成。1306 年 10 月以后：$10\times17.63=176.3$ 法郎，是 1305 年的 2.87 倍。贬值时，按记账货币收固定租金的房东与债权人受损，租户与债务人得益；恢复时反过来。德瓦伊所说的骚乱，正是租户在恢复时承受的那一下。契约写成名义数额，货币单位的价值一变，双方之间的财富就被重新分配了（\chref{ch:prices}）。
\end{example}

腓力四世之后，变动并没有停止。1342 年 7 月，里弗尔跌到 4.07 法郎。1351—1360 年，正值百年战争，里弗尔的价值变动了 71 次，其中 1359 年 16 次、1360 年 17 次；最高 13.59 法郎，最低 3.22 法郎，最高是最低的 4.2 倍，下跌之后又回升的情形有 11 次\cite{wailly1857}。

\begin{compare}{铸币税与通胀税}{中世纪的铸币税}{今天的通胀税}
\cmprow{怎样征收}{铸造时从送来的金属里扣下 $ME-MP$；贬值时扣得更多}{发行货币，使已有的货币贬值}
\cmprow{税基}{送进铸币厂的金银}{公众持有的货币}
\cmprow{谁承担}{送金属的人；按记账货币领取固定收入的人与债权人}{货币的持有者与名义债权人}
\cmprow{能收多少}{取决于有多少金属送来；贬值年份可占财政收入的一半以上}{取决于货币需求；通胀太高时人们少持货币，税基萎缩（\chref{ch:hyperinflation}）}
\cmprow{能否被察觉}{奥雷姆：比其他征收更难被察觉}{物价上涨看得见，税负落在谁身上却不容易看清}
\cmpsum{两者都是靠货币发行权取得的收入；贬值把正常年份微不足道的铸币税，变成了战时的主要财源。}
\end{compare}

\section{奥雷姆的批评}\label{sec:medieval-oresme}

这一节回答：当时的人怎样看待君主改变货币？法国神学家尼科尔·奥雷姆（Nicole Oresme，约 1320—1382）的《论货币》（\emph{De moneta}）是最系统的回答。奥雷姆在巴黎大学任教，为查理五世（Charles V，1364—1380 年在位）把亚里士多德的著作译成法文，晚年任利雪（Lisieux）主教。沃洛夫斯基判断，《论货币》“写成于 1373 年以前”；到 1373 年，奥雷姆自己的法译本已经藏在卢浮宫塔楼的王室图书馆里\cite{oresme1864}。沃洛夫斯基的校订本所据的抄本 A，抄写于 1397 年。

奥雷姆的出发点是：货币不是君主的财产。他在第 6 章说，“货币属于共同体，也属于个人”（\emph{Est igitur pecunia communitatis et singularium personarum}）\cite{oresme1864}。由此他得出了五个论点。下面的中文是本书的翻译。

\textbf{第一，君主的所得就是人民的所失。}第 15 章：“君主从中得到多少利润，共同体就必然受到多少损失。”这与今天说的\textbf{通胀税}是同一个道理：发行货币得到的收入不是凭空产生的，它是持币人手中购买力的转移（\chref{ch:budget}）。奥雷姆还说，从改变货币中取利，“比高利贷还坏”，因为在不正当的货币变更中，君主是“违背臣民的意愿，直接拿走他们的钱”（第 17 章）\cite{oresme1864}。上一章（\chref{ch:islamic}）讲过教会怎样禁止高利贷；奥雷姆用同一套道德语言，把矛头转向了君主。

\textbf{第二，一次变更会引出更多变更。}第 15 章：“如果君主有权作一次简单的货币变更并从中取利，同理，他也可以作更大的变更，取更大的利。”\cite{oresme1864}只要承认君主有权从改变货币中取利，就找不到一条界线阻止他下一次改得更多、取得更多。用今天的话说，这是\textbf{承诺}的问题：一个可以随时改变货币的君主，事先说的话不可信。高级教士在 1304 年前后请求腓力四世宣誓约束自己，被拒绝了；今天许多国家用央行独立来解决同一个问题（\chref{ch:credibility}）。

\textbf{第三，这是一种隐蔽的税。}第 20 章：“这样的征敛越不容易被察觉，施行起来就越危险；因为人民不像对其他征收那样很快感到它的负担。”\cite{oresme1864}议会税要召集议会、派人征收，人人看得见；贬值只需改一道铸币的命令。

\textbf{第四，金银会流出王国。}同一章：“由于这样的变更与贬损，王国里的金银减少了；因为尽管有人看守，它们还是被运到国外，在那里作价更高。”\cite{oresme1864}当本国把同样的金银作价得更低（或把劣币强行按足值作价）时，足值的金银就会流向估价更高的地方。这与后来所说的“劣币驱逐良币”是同一个道理。

\textbf{第五，消息灵通的人发财，固定收入的人受损。}第 21 章说，君主若不事先公布变更的时间与方式，“有些人靠心计或朋友暗中预知此事，于是用劣币买进货物，日后再按良币卖出，一下子就富了”；而且“按钱数定下的收入，必然不是被不公正地减少，就是被不公正地增加”\cite{oresme1864}。前一句说的是内幕消息，后一句说的正是例~\ref{ex:medieval-rent} 里的房东与租户。

奥雷姆没有统计数据，他的论证全靠推理，但这五点几乎覆盖了今天关于通货膨胀的主要议题：谁征了税，为什么政府管不住自己，为什么税负看不见，资本怎样外逃，以及名义契约怎样重新分配财富。

\section{数据检验与反例}\label{sec:medieval-test}

前面几节提出了几个论断：人口减少会提高工资、降低地租（\secref{sec:medieval-malthus}）；法令压不住工资（\secref{sec:medieval-statute}）；贬值是君主的财源（\secref{sec:medieval-debasement}）。这一节用数据逐一检验，看看哪些地方需要修正。

\subsubsection{方向对了，幅度不对}
图~\ref{fig:plague_wages} 把 Clark 的实际工资与实际地租，与式~\eqref{eq:medieval-factor} 的预测放在一起，都以 1300—1349 年的平均为 1。实际工资确实上涨了，实际地租也确实比瘟疫前低。可是幅度差得很远。1400 年代，模型预测实际工资为 1.30，实际的建筑小工是 1.83，农工 2.26，建筑工匠 1.45。模型预测实际地租为 0.68，实际是 0.97，几乎没有下降；1350—1460 年代平均约 0.91，最低的 1350 年代与 1370 年代也有 0.82，直到 1470 年代才降到 0.71。农工的工资不但远高于模型，而且高于 Clark 在工作论文里给出的整个范围：按不同的技术假设，瘟疫可以让工资“上涨 0 到 86\% 之间的任何幅度”\cite{clark1999wp}，而农工的实际工资在 1400 年代高出瘟疫前 126\%。

\fig{plague_wages}{1250—1500 年英格兰的人口、实际工资与实际地租。(a) 看两个人口序列都在 1348—1349 年骤降，此后一个世纪没有恢复；两者的水平不同。(b) 看实线与虚线的距离：实际工资涨得比模型预测的多得多，农工尤甚。(c) 看实际地租：模型预测下降三成，实际多在瘟疫前的九成上下，1470 年代以后才降到七成多。(b)(c) 以 1300—1349 年的平均为 1；虚线是按 Clark 的人口序列与式~\eqref{eq:medieval-factor}（$\alpha=0.405$）算出的预测。数据：Clark 整理的十年平均数（Clark 2010 说明）；实际地租为名义地租除以生活费用指数；Broadberry 等（2015），经英格兰银行 \emph{A Millennium of Macroeconomic Data}。}{fig:plague_wages}

\subsubsection{时间也不对}
模型说，劳动减少的那一刻工资就该上涨。农工的工资大致如此：1350 年代已是瘟疫前的 1.71 倍，而模型按那时的人口（瘟疫前的 0.71）只预测 1.15。建筑小工的工资却是逐步上升的：1350 年代 1.24，1360 年代 1.44，1370 年代 1.55，1380 年代 1.84，1440 年代才到 2.26。Clark 的工作论文用较早的农业工资数据，发现 1350—1369 年的实际工资甚至“略低于 1300—1349 年的平均”，并推测这“可能是《劳工法令》造成的假象”，因为法令把高于 1346—1347 年的工资定为非法\cite{clark1999wp}。在修订后的数据里，农工没有这样的滞后，建筑工人仍有；原因是什么，本章的材料不足以判断。

\subsubsection{效率提高了吗}
幅度的差异有一种简单的解释：瘟疫之后，同样的劳动与土地能生产更多的东西。Clark 给出了检验的办法。如果效率不变、资本回报率不变，工资的上涨与地租的下降应当大致相当，例~\ref{ex:medieval-plague} 里的“上涨 40\%、下降 39\%”就是一例。实际上工资上涨了 83\%（建筑小工），地租只下降了 3\%，两者都在模型之上。

Clark 衡量效率的公式是，效率的变化等于各要素实际价格变化的加权几何平均，权重是各要素的份额\cite{clark1999wp}：
\[
\frac{E_1}{E_0}=\Big(\frac{r_1}{r_0}\Big)^{0.16}\Big(\frac{w_1}{w_0}\Big)^{0.50}\Big(\frac{s_1}{s_0}\Big)^{0.34} .
\]
Clark 估计资本回报率从“1300 年前后的 10\% 降到 1400 年的约 5\%”；按 1300—1348 年永久租费（rent charge）的买卖价格算，中位回报率为 8.1\%\cite{clark1999wp}。取 $r_1/r_0=5/8.1\approx0.62$，用 1400 年代的实际工资与地租，建筑小工给出 $E_1/E_0\approx1.24$，农工给出 1.38。如果工资与地租恰好按模型变化、回报率不变，这个比值正好是 1。这是本章用 Clark 2015 年的数据和公式做的粗略计算，不是 Clark 发表的结果；我们用生活费用指数而不是农产品价格去平减，这会影响结果。

Clark 自己的结论前后不同。1999 年的工作论文说，数据“没有任何迹象表明黑死病对农业的效率有任何影响”\cite{clark1999wp}；那篇论文的实际工资“是用农产品价格平减的名义工资”，15 世纪的实际农业工资只比瘟疫前高 64\%。2016 年正式发表的论文摘要则说，黑死病的人口损失“伴随着经济效率的意外提高，尽管经济的规模缩小了；但当人口在 16 世纪再次增长时，这种效率的提高消失了”\cite{clark2016}。同一位作者、同一个问题，换了数据与平减指数，结论就变了。

\begin{pitfall}[日工资不是生活水平]
图~\ref{fig:plague_wages} 的实际工资是\textbf{日工资}除以生活费用指数。它有两个陷阱。第一，用哪一套工资数据，结论很不一样：Clark 的工作论文用 Farmer 整理的计件工资折算日工资，再用农产品价格平减，15 世纪的实际农业工资高出瘟疫前 64\%\cite{clark1999wp}；Clark 后来整理的新序列里，15 世纪农工的平均实际工资高出 141\%。差别主要不在平减指数：新序列的生活费用指数在 15 世纪平均只比瘟疫前低 1\%，141\% 几乎全部来自名义日工资的上涨（平均为瘟疫前的 2.44 倍）；工作论文的农产品价格指数在 15 世纪平均也与瘟疫前相差无几，它的 64\% 之所以低得多，是因为由计件工资折算出的名义日工资涨得少得多\cite{clark1999wp}。第二，日工资乘以工作天数才是年收入。Humphries 与 Weisdorf 用按年雇用的工人的收入推算，按日计酬的工人要挣到同样的年收入，中世纪一年只需工作“不到 150 天”，工业革命时期则要“远超 300 天”\cite{hw2019}。他们的新估计“下调了”黑死病之后的“劳动者的黄金时代”：即使在那时，一个没有技能的男工“也很难养活自己以外的一个人以上”\cite{hw2019}。日工资翻倍，不等于生活水平翻倍；还要问一年能干多少天、愿意干多少天。
\end{pitfall}

\subsubsection{人口为什么没有恢复}
模型的第三个预测是人口回升：取 $\alpha=0.405$，增幅的一半在 1.3 代之内消失。可是按 Broadberry 等人的序列，1450 年的人口只有 1348 年的 40\%，比 1351 年还少；Clark 的序列也显示，人口直到 15 世纪末仍远低于瘟疫前（图~\ref{fig:plague_wages}(a)）。工资之所以长期居高不下，首先是因为劳动一直稀缺。按\secref{sec:medieval-malthus}末尾的解读，反复的瘟疫提高了养活一个成年后代的成本 $\rho$，稳态收入随之提高；Clark 2016 年的摘要说效率的提高在 16 世纪人口增长时消失，也与此相容。人口长期没有恢复，并不否定马尔萨斯模型；它说明不能把黑死病当成一次性的冲击。

\subsubsection{贬值之谜}
\secref{sec:medieval-debasement}的数据支持“贬值是财源”的说法：贬值的年份，铸币收入占财政收入的一半以上。但是，人们为什么愿意在贬值时把金属送进铸币厂？Rolnick、Velde 与 Weber 指出：“如果钱币在流通中按其内含的金属估值，贬值就不可能给人们任何额外的动力把金属送进铸币厂。可是贬值确实吸引了大量的金属”；加强之后的铸币量“也异常地大，与贬值之后差不多”\cite{rvw1996}。按他们的说法，“就我们所知，即使在贬值期间，金属也是自愿送到铸币厂的”\cite{rvw1996}。

流行的解释是钱币按枚数（by tale）而不是按所含的金属流通：新旧钱币按同一面值通用，拿足重的旧币去改铸成轻的新币就有利可图。Rolnick 等用当时的证据反驳：贬值之后新旧钱币同时流通，至少有时按各自所含的金属比价交换。这种解释的一个变体是价格调整得慢，新币在一段时间里还能按旧的物价花出去；他们认为，就算价格不是立即调整，“滞后也只能是几个星期”。还有一种解释是贬值减轻了按名义数额计的债务，他们同样不接受\cite{rvw1996}。这与奥雷姆的第五点和巴黎租户的骚乱并不矛盾：名义契约确实被重新分配了，只是这不足以解释为什么那么多金属被送进铸币厂。他们的结论是：“否定了现有的解释，我们面临的就是一个建模的挑战：解开贬值之谜”\cite{rvw1996}。本章也把这个问题留作开放的问题。

还有一个问题是，英格兰为什么比法国稳定得多。英镑在 1542 年之前一次从未损失超过 20\%，法国却有 123 次贬值。一个线索是战争。Rolnick、Velde 与 Weber 估计，与 1330 年代或 1370 年代的财政收入相比，1350 年代法国铸币税占财政收入的比例“在 8\% 到 12\% 之间”，并且“十有八九实际高得多，因为国家正处在战争中，正常的税收大概处于低谷”\cite{rvw1996}。打仗需要钱，改铸币又不必召集议会，比征税快得多；而按奥雷姆的第三点，它的负担也不像其他征收那样很快被人感到。至于英格兰为什么较少走这条路，本章的材料不足以回答。

\begin{compare}{实物冲击与货币冲击}{黑死病（实物冲击）}{铸币贬值（货币冲击）}
\cmprow{改变了什么}{劳动与土地的相对数量}{记账单位所含的金银}
\cmprow{相对价格}{实际工资上升、实际地租下降}{名义价格变动；按记账货币写的契约的实际价值变动}
\cmprow{谁得谁失}{幸存的劳动者得益，地主受损}{君主取得铸币税；贬值时债务人与租户得益，恢复时反过来}
\cmprow{持续多久}{几代人：取决于人口何时恢复}{价格调整只需“几个星期”，但对既有契约的再分配已经发生}
\cmprow{今天的对应}{技术与供给冲击（\chref{ch:rbc}）；人口变化（\chref{ch:demography}）}{通货膨胀与货币政策（\chref{ch:prices}、\chref{ch:hyperinflation}）}
\cmpsum{实物冲击改变经济能生产什么、谁的投入更稀缺；货币冲击改变计价的尺度，它的实际后果主要来自没有随尺度调整的契约。}
\end{compare}

\section{今天还剩什么}\label{sec:medieval-legacy}

庄园与贬值的里弗尔都已成为历史，但这一章的几个制度要素与分析框架今天还在。

\textbf{马尔萨斯的约束与它的解除。}式~\eqref{eq:medieval-dynamics} 说，在土地固定、技术停滞的经济里，人均收入长期只取决于生育与养育的成本。14 世纪的英格兰离开稳态，靠的是一场灾难；后来的欧洲怎样持续地摆脱这个约束，是\chref{ch:divergence}的问题，而技术进步怎样代替土地成为增长的来源，见\chref{ch:solow}。今天许多国家面对的是相反的问题：人口减少而不是增加，劳动变得稀缺，这对工资与增长意味着什么，\chref{ch:demography}会回到这个问题。

\textbf{管制价格与劳动市场。}《劳工法令》把工资压在旧水平，结果是短缺、规避与对流动的限制。这与\chref{ch:greekroman}的戴克里先限价令、后来各种价格管制的逻辑相同；劳动市场的供求与工资刚性，见\chref{ch:labor}。

\textbf{铸币税与通胀税。}中世纪君主靠贬值取得的收入，今天以发行货币的形式存在。它在政府预算中的地位见\chref{ch:budget}；当它被推到极端，就是\chref{ch:hyperinflation}的恶性通胀；当央行被迫为财政融资，就是\chref{ch:fiscaldominance}的财政主导。

\textbf{承诺与可信度。}高级教士请求腓力四世宣誓只在“绝对必要”时才改变货币，被拒绝了；奥雷姆说，允许一次贬值就挡不住下一次。今天的答案是把货币交给独立的央行，并公开承诺通胀目标（\chref{ch:credibility}）。

\textbf{名义契约。}巴黎的租金写的是里弗尔，里弗尔的金银含量一变，房东与租户之间的财富就被重新分配。今天的债务、工资与租金大多也写成名义数额，通货膨胀以同样的方式重新分配财富，这是\chref{ch:prices}区分名义与实际的原因。16 世纪的价格革命究竟来自人口还是美洲白银，\chref{ch:silver}会把本章的两类冲击放在一起比较。

\begin{deeper}[单位、数据与几处不一致]
\textbf{单位。}英格兰的记账单位是 1 镑 = 20 先令 = 240 便士；文献里的 l.、s.、d. 分别指镑、先令、便士，ob.（obolus）指半便士。法国是 1 里弗尔 = 20 苏 = 240 德尼耶；本章的里弗尔都指图尔里弗尔（livre tournois）。德瓦伊的法郎是 1857 年的法郎，一法郎为 5 克九成成色的银；他的表格把里弗尔的价值算到一个生丁以下七位小数，本章照他正文的写法截取到生丁。马克（mark）是贵金属的重量单位。

\textbf{两个人口序列。}Broadberry 等人的序列与 Clark 的序列水平不同：1340 年代 Clark 为 448 万，Broadberry 等 1340 年为 456 万、1348 年为 481 万；1400 年代 Clark 为 264 万，Broadberry 等 1400 年为 208 万。本章凡与 Clark 的工资对比，都用 Clark 的人口；单独描述人口时用 Broadberry 等人的序列，并注明来源。实际地租是 Clark 的名义地租除以其生活费用指数，只能作为指数，不能读作镑数。

\textbf{几处不一致。}（1）南弗里斯 1349—1350 年的账目逐项相加为 118 先令 10¼ 便士，比账上的合计少 5 便士；本章只用账上写明的合计。（2）1349 年的条例以第二十年“或此前五六个普通年份”的工资为准，1351 年的法令对农业仆役写的是“或此前四年”。（3）Rolnick、Velde 与 Weber 的正文说大贬值期间铸币税占英格兰王室收入的 25\%，他们的表 3 是 26\%（1544—1547 年），本章用表中的数字。（4）Clark 1999 年的工作论文认为黑死病没有提高效率，2016 年的论文认为暂时提高了，本章两者都引，并注明各是哪一篇。（5）网上流传的一份题为“1351 年劳工法令”的英译（耶鲁大学 Avalon 项目），内容其实是 1349 年的条例；本章不用它，条例据 Bland、Brown 与 Tawney 的译文，法令据 \emph{Statutes of the Realm} 法文原文旁的英译。
\end{deeper}

\begin{takeaways}
\begin{itemize}
\item 黑死病（1348—1350）使英格兰人口在三年里减少约 46\%，而且一个世纪没有恢复；土地没有减少，劳动变得稀缺。
\item 马尔萨斯模型预测，劳动减少使工资上升、地租下降，幅度由土地的份额决定：取 $\alpha=0.405$，人口降到 0.53，工资上涨 30\%、地租下降 32\%，增幅在几代之内消失。
\item 数据在方向上支持模型，在幅度上不支持：1400 年代的实际日工资是瘟疫前的 1.45—2.26 倍，实际地租几乎没有下降。效率的提高、反复的瘟疫与日工资本身的缺陷，都是可能的解释。
\item 1349 年的条例与 1351 年的法令把工资压在瘟疫前的水平，形成了一个价格上限下的短缺；严厉的执行（温斯特里百户区九成的税款来自工人的罚款）也没能阻止工资上涨，立法对劳动者流动的限制也越收越紧。
\item 铸币贬值是君主的财源：法国贬值的年份，铸币收入占财政收入的一半以上，正常年份通常在 5\% 以下。奥雷姆指出，它是隐蔽的征税，会自我加码，使金银外流，让知情者获利，并重新分配名义契约；为什么贬值时人们仍把金属送进铸币厂，至今仍是一个谜。
\end{itemize}
\end{takeaways}

\begin{selfcheck}
\begin{enumerate}
\item 哈弗林的清册把每个维尔盖特的劳役折价为每年 7 先令。瘟疫之后雇工的工资大涨，而劳役的折价不变。领主更愿意要劳役还是要折算的钱？佃户呢？
\item 设劳动降到瘟疫前的 0.6，$\alpha=0.405$。实际工资与地租各变化多少？如果此后不再有瘟疫，一代之后人均收入比稳态高多少？
\item 1350 年的令状规定，工人超额所得的工资要交给所属村镇，用来缴纳议会税。这对村民举报多拿工资的人有什么影响？用温斯特里百户区的数字说明。
\item 设某年铸币厂的 $MP/ME=85\%$，铸造成本为 3\%。毛铸币税率与净铸币税率各是多少？如果送来的金属是正常年份（毛铸币税率 7.5\%）的 3 倍，净铸币税收入是正常年份的几倍？
\item 在例~\ref{ex:medieval-rent} 中，10 里弗尔的租金在 1303 年 8 月合多少法郎？是 1295 年以前的几成？
\item 用 Clark 的效率公式：如果效率与资本回报率都不变，而实际工资上涨到原来的 1.83 倍，实际地租应当变成原来的多少？实际是 0.97，这说明了什么？
\end{enumerate}
\end{selfcheck}

\begin{answers}
\item 雇工的工资涨了，而劳役仍按旧价折算，所以对领主来说，劳役比收钱再雇人划算，他更想要劳役；对佃户来说，同样的劳动拿到市场上能换更多的钱，他更想把劳役折成钱，按旧价交钱就可以脱身。双方的利益在瘟疫后变得对立。
\item 由式~\eqref{eq:medieval-factor}，$0.6^{-0.405}\approx1.23$，$0.6^{0.595}\approx0.74$：工资上涨 23\%，地租下降 26\%。一代之后，$1.23^{0.595}\approx1.13$，人均收入比稳态高 13\%。
\item 超额的工资变成了本村的税款，举报一个多拿工资的人，就减轻了全村的税负，村民（尤其是雇主）有了举报与执行的动力。温斯特里百户区 1351—1352 年的税款中，工人罚款占 90\%，一半村镇的税额完全由罚款抵付。但执行越严，越说明多付工资的交易有多普遍。
\item 毛铸币税率 $1-85\%=15\%$，净铸币税率 $15\%-3\%=12\%$。正常年份的净铸币税率为 $7.5\%-3\%=4.5\%$；净收入之比为 $3\times12/4.5=8.0$ 倍。
\item $10\times6.97=69.7$ 法郎，是 1295 年以前 201.0 法郎的 $69.7/201.0\approx34.7\%$，不到三分之一。这与德瓦伊说的租金缩水到“不足原值的三分之一”大体一致。
\item 效率与回报率不变时，$1.83^{0.50}(s_1/s_0)^{0.34}=1$，得 $s_1/s_0=1.83^{-0.50/0.34}\approx0.41$：地租应当下降 59\%。实际只下降 3\%，说明工资与地租不可能同时都按“效率不变”的方式变化：或者效率提高了，或者工资与地租的数据（日工资、平减指数）有偏差，或者两者兼有。
\end{answers}

欧洲的君主靠改变铸币的金银含量取得收入。同一时期的中国，宋、元两代已经在发行与金属脱钩的纸币，财政压力也以另一种方式变成了通货膨胀。下一章转向东方，看世界上最早的纸币怎样兴起，又怎样崩溃。
